% answers_short.tex —— 2025 年·全国 2 卷 · 简短答案速查
% 由 tools/make_short_answers.py 自 answers.tex 的 \gksolhead/\gkanswer 提取，
% 并与 2025/2025_数学真题_全国2_答案要点.md 逐题比对后生成，勿手改（改源 MD 后重跑）。

\begin{gknotice}
\textbf{使用说明}：本节只给最终答案便于快速核对，完整推导见解析版。答案为编写者独立推导，\textbf{非官方标准答案}。
\end{gknotice}

\gkpart{单选题}
\begin{center}\begin{tabular}{ccccccccc}
\toprule
题号 & 1 & 2 & 3 & 4 & 5 & 6 & 7 & 8 \\
\midrule
答案 & C & A & D & C & A & C & B & D \\
\bottomrule
\end{tabular}\end{center}

\gkpart{多选题}
\begin{center}\begin{tabular}{cccc}
\toprule
题号 & 9 & 10 & 11 \\
\midrule
答案 & AD & ABD & ACD \\
\bottomrule
\end{tabular}\end{center}

\gkpart{填空题}
\begin{itemize}[leftmargin=2.6em]
  \item[\textbf{12.}] $\sqrt2$
  \item[\textbf{13.}] $-4$
  \item[\textbf{14.}] $\dfrac52$
\end{itemize}

\gkpart{解答题}
\begin{itemize}[leftmargin=2.6em]
  \item[\textbf{15.}] (1) $\varphi=\dfrac\pi3$；(2) 值域 $\left[-\sqrt3,\sqrt3\right]$；递减区间 $\left[k\pi-\dfrac\pi{12},\ k\pi+\dfrac{5\pi}{12}\right]$，递增区间 $\left[k\pi+\dfrac{5\pi}{12},\ k\pi+\dfrac{11\pi}{12}\right]$（$k\in\mathbf Z$）
  \item[\textbf{16.}] (1) $\dfrac{x^2}{4}+\dfrac{y^2}{2}=1$；(2) $|AB|=\sqrt5$
  \item[\textbf{17.}] (1) 见解析版；(2) $\dfrac{\sqrt{42}}{7}$
  \item[\textbf{18.}] (1) 见解析版；(2) ② $2x_1>x_2$
  \item[\textbf{19.}] (1) $p_3=p^3$，$p_4=4p^3-3p^4$；(2) $p=\dfrac23$；(3) 见解析版
\end{itemize}
